\subsection{Arbitrary two-by-three minimal bands opposite a right-zero corner}

\begin{theorem}[Proved: every two-by-three minimal band opposite $R_3$]
\label{CP2:A-two-by-three-band-general-right-zero}
Let $M$ be any finite monoid with proper minimal ideal a
$2\times3$ rectangular band.  For every finite left-$R_3$,
right-$M$ biset, all rational degree-two products of the
associated collage vanish.  No bound on the source order,
state count, homology, or second Betti number is required.
\end{theorem}
\begin{proof}
Write $K=\{e_{ij}:i\in\{0,1\},\ j\in\{0,1,2\}\}$,
with $e_{ij}e_{uv}=e_{iv}$, and write
\[
 me_{ij}=e_{\sigma_m(i),j},\qquad
 e_{ij}m=e_{i,\tau_m(j)}.
\]
Let $\chi(m)\in\{0,1,-1\}$ be the action of $\sigma_m$
on the one-dimensional row-augmentation module.
Let $W=\ker(\mathbb Q\{j_0,j_1,j_2\}\to\mathbb Q)$.
The natural balanced presentation of Lemma~\ref{A:band-Tor}
is
\[
 H_2(BM;\mathbb Q)
   =W/\langle w\tau_m-\chi(m)w:m\in M,\ w\in W\rangle.
                                                               \tag{1}
\]
Its dimension $d$ is zero, one, or two.
If $d=0$, both corners have zero rational second cohomology.
All second classes of the collage are then relative to the
two corners, and their products vanish by the actual
suspension quotient in the two-block model.
Assume $d>0$.

First suppose $d=2$.  Every relation in (1) is zero, so
the induced action of $\tau_m$ on $W$ is $\chi(m)I$.
If $\chi(m)=0$, this forces $\tau_m$ to be constant.
If $\chi(m)=1$, the action has rank two, so $\tau_m$
is a permutation, and its identity action on the augmentation
space forces it to be the identity permutation.
The alternative $\chi(m)=-1$ is impossible: a permutation
of three elements acting by $-I$ on $W$ would have trace
$1-2=-1$ on the full three-dimensional permutation module,
whereas that trace is its nonnegative number of fixed points.
Thus the simultaneous actor image consists of the six
constant actor pairs and the identity pair, exactly the
identity-adjoined $2\times3$ band $P_3$.
The actual actor quotient $M\twoheadrightarrow P_3$
identifies the band bijectively.  By
Proposition~\ref{CP2:A-band-quotient-collage-H2}, it induces
an isomorphism in the second cohomology of collages after
taking the actual induced state quotient.  The target has
zero degree-two products by
Theorem~\ref{CP2:A-two-by-three-band-right-zero-all-states}.
Their actual pullbacks prove the assertion in this case.

It remains to consider $d=1$.  Choose a nonzero linear
functional $\lambda$ on the quotient (1), and put
\[
 a_0=0,\qquad a_j=\lambda(j_j-j_0)\quad(j=1,2).
\]
The three rational values $a_j$ are not all equal.
The balancing equations say, literally,
\[
 a_{\tau_m(j)}=\chi(m)a_j+t_m,
 \qquad t_m=a_{\tau_m(0)}.                                  \tag{2}
\]
Consequently every column actor induces an affine map
of the finite value set $T=\{a_0,a_1,a_2\}$, and equal
values are preserved by every column actor.

Suppose $T$ has two elements.  Let $\pi:\{0,1,2\}\to T$
send $j$ to $a_j$, and let $\overline\tau_m$ be the induced
transformation of $T$.  Equation (2) shows that its
augmentation scalar is exactly $\chi(m)$.
Indeed a zero slope is constant; slope $1$ is a translation
of a finite set and is therefore the identity; and slope
$-1$ interchanges the two values.
The homomorphism
\[
 \rho:M\longrightarrow H,
 \qquad m\longmapsto(\sigma_m,\overline\tau_m)
                                                               \tag{3}
\]
has as its image $P_5$ or the six-element simultaneous-swap
extension of $P_5$.  All four constant pairs occur, because
the band maps onto the two rows and two value columns.

Although (3) identifies two original band columns, it
identifies the minimal left ideal $Me_{00}$ bijectively
with $H\rho(e_{00})$: both have the same two row labels,
and their left $M$ actions agree.  Thus
Lemma~\ref{CP2:A-idempotent-quotient-bridge-connectivity}
applies to the actual quotient
$U=S\otimes_M H$ and gives a bridge $H_1$ isomorphism
and a bridge $H_2$ surjection, together with the corner
$H_1$ isomorphism.

The corner $H_2$ map is also an isomorphism.
Both its source and target have dimension one by (1)
and the equal-character presentation of the two-column
target.  It is nonzero: the functional on the target
column augmentation assigning the two values in $T$
pulls back along $\pi$ to the chosen nonzero $\lambda$.
This is the actual natural map of the two augmentation
sequences, so the nonzero map is the one induced by (3).
In the natural two-block exact sequence, the bridge
$H_2$ surjection identifies the preceding cokernels,
and the bridge and corner $H_1$ isomorphisms identify
the following kernels.  Hence the actual collage map
is an $H_2$ isomorphism, as in the proof of
Proposition~\ref{CP2:A-band-quotient-collage-H2}.
Its cohomological pullback is therefore an isomorphism.
The target is a two-by-two band collage and has zero
degree-two products by
Theorem~\ref{CP2:A-two-by-two-band-right-zero-all-states}.
All source products vanish as their actual pullbacks.

Finally suppose $T$ has three elements, so the column
labels are distinguished by their values.  Equation (2)
forces a zero-slope actor to be constant and a slope-one
actor to be the identity, again because a finite subset
of $\mathbb Q$ admits no nonzero translation preserving it.
A slope-minus-one actor must be a reflection of $T$.
It interchanges its minimum and maximum and fixes its
middle value; it can exist only when the middle value
is their arithmetic mean.  This is the unique possible
reflection, and its simultaneous row actor is the swap.
Thus the faithful simultaneous actor image is either
$P_3$ or its eight-element extension by this simultaneous
row swap and column reflection.
The actor map is bijective on the original six band
elements, so Proposition~\ref{CP2:A-band-quotient-collage-H2}
again gives an actual second-cohomology isomorphism of
collages.  Relabeling the middle column as $0$ identifies
the eight-element image with source $500$ of
Theorem~\ref{CP2:A-eight-two-by-three-all-sources}.
That theorem proves its all-state product vanishing;
the pure-band image is covered by the corresponding
all-state pure-band theorem.  Pullback completes this case.

If the connecting set is empty, the same corner maps
are $H_2$ isomorphisms and the collage is a disjoint
union.  Each image monoid above has trivial rational
module of projective dimension at most two: this is
the pure-band resolution, the semisimple contracted
two-column resolution, or the displayed complete
source-$500$ resolution.  Its fourth cohomology is
therefore zero, and the same actual pullback argument
proves the assertion for an empty set as well.
\end{proof}
